\documentclass[11pt,a4paper,oneside]{article}

% --- Поддержка русского языка в pdfLaTeX ---
\usepackage[T2A]{fontenc}
\usepackage[utf8]{inputenc}
\usepackage[russian]{babel}

% --- Математические библиотеки ---
\usepackage{amsmath,amssymb,amsfonts,amsthm,mathtools}

% --- Геометрия страницы и типографика ---
\usepackage[left=20mm, right=20mm, top=22mm, bottom=22mm, headheight=14pt]{geometry}
\usepackage{microtype}
\usepackage{tabularx}
\usepackage{array}
\usepackage{booktabs}
\usepackage{verbatim}
\usepackage{xcolor}
\usepackage{fancyhdr}
\usepackage{hyperref}

\hypersetup{
    colorlinks=true,
    linkcolor=blue!70!black,
    urlcolor=blue!70!black,
    citecolor=blue!70!black,
    pdftitle={Задача A. Тепловые панели},
    pdfauthor={Яндекс CodeRun}
}



% --- Настройка олимпиадных макросов Polygon / Codeforces ---
\newenvironment{problem}[5]{
    \addcontentsline{toc}{section}{#1}
    \begin{center}
        {\LARGE\bfseries #1}\par
        \vspace{0.3cm}
        {\small\sffamily
        \begin{tabular}{rl}
            \textbf{Входной файл:} & #2 \\
            \textbf{Выходной файл:} & #3 \\
            \textbf{Ограничение по времени:} & #4 \\
            \textbf{Ограничение по памяти:} & #5 \\
        \end{tabular}
        }\par
    \end{center}
    \vspace{0.3cm}
    \par\noindent
}{
    \vspace{0.8cm}
}

% --- Секции олимпиадного условия ---
\newcommand{\InputFile}{\subsubsection*{Формат входных данных}}
\newcommand{\OutputFile}{\subsubsection*{Формат выходных данных}}
\newcommand{\Examples}{\subsubsection*{Примеры}}
\newcommand{\Example}{\subsubsection*{Пример}}
\newcommand{\Note}{\subsubsection*{Примечания}}

% --- Колонтитулы ---
\pagestyle{fancy}
\fancyhf{}
\fancyhead[L]{\small\sffamily Задача A. Тепловые панели}
\fancyhead[R]{\small\sffamily Яндекс CodeRun (2025-winter-common)}
\fancyfoot[C]{\thepage}
\renewcommand{\headrulewidth}{0.4pt}
\renewcommand{\footrulewidth}{0pt}

\setlength{\parindent}{1.5em}
\setlength{\parskip}{0.4ex plus 0.2ex minus 0.1ex}

\begin{document}

% ----------------------------------------------------------------------
% Задача A. Тепловые панели
% Тема: Теория чисел и геометрия прямоугольников | День: 1
% ----------------------------------------------------------------------
\begin{problem}{A. Тепловые панели}{стандартный ввод}{стандартный вывод}{3.0 секунда}{512 МБ}

Лорд Нуль заморозил Солнечный процессор, и~на цифровой мир опустилась Вечная зима. Становится всё холоднее. В~тренировочном лагере у~подножья Пика Кода Кодерун учится собирать технические панели для обогрева. Каждая панель состоит из~красных плиток-контуров $R$, которые формируют защитную рамку, и~синих плиток-ядра $B$, аккумулирующих тепло внутри.

Нужно определить размеры такой прямоугольной панели~---~ширину $W$ и~высоту $H$ ($W \ge H$)~---~зная количество красных и~синих плиток, при~условии, что красные плитки находятся по~периметру стены в~один слой, а~синие~---~внутри.

\InputFile
В~единственной строке вводятся два числа: $R$ ($8 \le R \le 10^9$) и~$B$ ($1 \le B \le 10^9$)~---~количество красных плиток по~периметру и~количество синих плиток внутри соответственно.

\OutputFile
Выведите размеры панели~---~два числа $W$ и~$H$, такое что $W \ge H$. Гарантируется, что ответ существует и~единственный.

\Examples
\begin{center}
\begin{tabular}{|p{0.48\textwidth}|p{0.48\textwidth}|}
\hline
\textbf{Стандартный ввод} & \textbf{Стандартный вывод} \\
\hline
\begin{minipage}[t]{0.47\textwidth}
\vspace{0pt}
\begin{verbatim}
10 2
\end{verbatim}
\vspace{0pt}
\end{minipage}
&
\begin{minipage}[t]{0.47\textwidth}
\vspace{0pt}
\begin{verbatim}
4 3
\end{verbatim}
\vspace{0pt}
\end{minipage} \\
\hline
\multicolumn{2}{|p{0.96\textwidth}|}{\small\textit{Пояснение к примеру 1:} Панель имеет размеры $4 \times 3$. Защитная рамка периметра содержит $2 \cdot 4 + 2 \cdot 3 - 4 = 10$ красных плиток, а~внутреннее ядро~---~$(4 - 2)(3 - 2) = 2$ синие плитки. Условие $W \ge H$ ($4 \ge 3$) выполнено.} \\
\hline
\begin{minipage}[t]{0.47\textwidth}
\vspace{0pt}
\begin{verbatim}
24 24
\end{verbatim}
\vspace{0pt}
\end{minipage}
&
\begin{minipage}[t]{0.47\textwidth}
\vspace{0pt}
\begin{verbatim}
8 6
\end{verbatim}
\vspace{0pt}
\end{minipage} \\
\hline
\multicolumn{2}{|p{0.96\textwidth}|}{\small\textit{Пояснение к примеру 2:} Панель имеет размеры $8 \times 6$. Защитная рамка периметра содержит $2 \cdot 8 + 2 \cdot 6 - 4 = 24$ красные плитки, а~внутреннее ядро~---~$(8 - 2)(6 - 2) = 24$ синие плитки. Условие $W \ge H$ ($8 \ge 6$) выполнено.} \\
\hline
\end{tabular}
\end{center}

\Note
Ответ всегда существует и~однозначно определяется геометрической системой уравнений: $2W + 2H - 4 = R$ и~$(W - 2)(H - 2) = B$ при~$W \ge H$.

\end{problem}


\end{document}
